\section{Measurement}\label{sec:measure}

The repository keeps four benches under \texttt{tools/}. Each is built to block one failure mode, and each has a scorer that is a
script. The numbers in this section come only from result files that are committed in the repository; each table says which. Where
a bench has no committed result for a claim, the text says so.

\subsection{Use-case bench: was the text read as a unit?}

\texttt{tools/usecase-bench} renders 52 scenes in six use cases, drawn from HTML written for the bench, with the scene
dimensions in \texttt{set/index.tsv}: 5 chats, 11 documents and web pages, 8 game screens, 7 papers, 9 subtitle frames and 12 code
scenes. For each scene a file lists the \emph{units} that a reader should return as one block (a paragraph, a table cell, a
title): 400 units in all (36 in chats, 89 in documents, 38 in games, 70 in papers, 20 in subtitles, 147 in code). A rendered scene is read, and
every unit gets one of five outcomes:
\emph{correct} (the unit is exactly one block), \emph{merged} (it is inside a longer block), \emph{split} (it is cut into pieces that
together cover at least $80\%$ of it), \emph{misread} (the best block is at least $0.75$ similar) and \emph{missing}. Eighty of the units are marked
as \emph{keep}, which must come back as written, and 21 carry inline substrings that must survive a translation verbatim; for these the scorer adds
three columns, \texttt{kept\_ok}, \texttt{translated\_ok} and \texttt{inline\_ok}. The scenes with odd numbers are the ones the paragraph-reader decision was tuned
on; the even ones are held out. The bench runs the three readers of Section~\ref{sec:overview}: line grouping, the document model, and the
automatic choice between them.

\subsection{Paper bench: does a typeset page survive?}

Typeset papers are a hard page: two columns, dense inline mathematics, tables, footnotes, a title block
with several columns of authors. \texttt{tools/paper-bench} renders every page of a paper PDF at a scale of $2.67$, as a Retina
screen showing a page about 820 points wide would, cuts two lens-sized crops per page whose top and bottom edges run through the middle of a line of text, and
cuts two pages with side crops that cut words at the left and right edge. Each render is scored by the metrics of
Table~\ref{tab:paper}. From this revision on, the page source of the bench is this paper: a document whose formulas, listings and tables
are ours to show.

\begin{table*}[t]
\centering
\caption{Metrics of the paper bench, from the docstring of \texttt{tools/paper-bench/score.py}.}\label{tab:paper}
\small
\begin{tabular}{@{}p{1.25in}p{5.5in}@{}}
\toprule
Metric & What it counts\\
\midrule
\code{formula\_translated} & translated sources that hold a sign of mathematics and are not a body sentence (fewer than six prose words)\\
\code{math\_heavy\_kept} & blocks left as written because they are mostly inline mathematics ($H$ of Eq.~\ref{eq:share})\\
\code{overlaps, spill, cut, clipped} & sums of the render's layout line; \emph{clipped}: a translation that needs more room than its patch has at the smallest size\\
\code{ellipsis} & of those, the translations cut to fit and ended in an ellipsis\\
\code{big\_text} & drawn body-like blocks (at least 40 characters per source line) whose font size is over $1.3$ times the median of the page's body-like blocks\\
\code{cross\_col} & pairs of blocks where one block's source rectangle holds at least half of the other's: a join of lines across a column gap\\
\code{debris} & translated paragraphs whose translation carries three or more one-letter tokens or a stray symbol that the OCR made of a typeset sign\\
\code{edge\_cut}, \code{side\_cut} & translated blocks whose source lines touch the top or bottom edge of a crop (within 4 px), or the left and right edge\\
\bottomrule
\end{tabular}
\end{table*}

The results of running the bench on the PDF of this paper are written beside the tool (\texttt{RESULTS-own-2026-10-19.txt}). They are
not repeated here: a number about this page would change the page.

\subsection{Keep-label bench: where do the rules put the boundary?}

\texttt{tools/keep-labels} asks the question of Section~\ref{sec:keep} directly: for every word that the reader returns, must it be
translated, or must it stay as written? Each word gets a label (\emph{keep}, \emph{translate}, \emph{either}) and a reason (mathematics, code, command, flag, path,
URL, identifier, hash, hotkey, log, configuration, line number, diff mark) from the typeset source: the mathematics of the PDF's
LaTeX source and the typewriter text of its fonts for papers, and the HTML for the code cases. Papers and cases each have a development set and a held-out set, and every set is read two ways: with the document model and with line grouping. The held-out paper is a physics paper under a
CC BY licence \citep{holdout}; the development papers are not shown here.

Table~\ref{tab:keep} gives the block-level result on the held-out sets with the document reader. A block is a true keep when its truth is keep words and no translate
word. Adding \code{CodeRegion} to the rules of one block raises the recall of the held-out code cases from $0.78$ to $1.00$ at a precision of $1.00$; on the
development cases recall goes from $0.56$ to $0.91$ and precision from $0.95$ to $0.94$. It does not change a single block of the held-out paper: the region
rule has nothing to say about prose with mathematics, and the one-block rules keep $26$ of its $44$ true keep blocks at three false positives. Of the $44$
case blocks whose truth is pure translation, the region rule fired on one: a line of an algorithm figure that is mostly symbols with a prose label.

\begin{table}[t]
\centering
\caption{Block-level keep precision and recall (document reader). Before: the rules of one block (\code{TranslationFilter}). After:
the same rules or a \code{CodeRegion} verdict. Source: \texttt{RESULTS-code-region-2026-10-19.txt}.}\label{tab:keep}
\small
\begin{tabular}{@{}lrrrrr@{}}
\toprule
Set & & P & R & tp/fp/fn\\
\midrule
Cases, development & before & 0.95 & 0.56 & 19/1/15\\
 & after & 0.94 & 0.91 & 31/2/3\\
Cases, held out & before & 1.00 & 0.78 & 31/0/9\\
 & after & 1.00 & 1.00 & 40/0/0\\
Paper, held out & before & 0.90 & 0.59 & 26/3/18\\
 & after & 0.90 & 0.59 & 26/3/18\\
\bottomrule
\end{tabular}
\end{table}

\subsection{Translation bench: what does context decide?}

\texttt{tools/translation-bench} measures what context does to the translation of screen text. Its corpus is written for the bench: 14 sequences
(5 chats, 3 documents, 3 games, 3 subtitles) with 445 lines, played back screen by screen the way a lens would see them, in batches of at most three new lines. In it
312 \emph{probes} are planted: a name introduced and reused, a dropped subject, a Korean speech level, a term that must keep one rendering, a domain sense. A probe records where
the fact that decides it lives: in the user memory only (126 probes), in earlier lines only (133), or nowhere, a control that any engine
should get right (53). Every probe is checked by a rule on the target line that must pass on a skilled human translation and fail on the untranslated source, so an
oracle that replays the references scores 312 of 312 under each of the seven conditions it was run with. The bench also defines a blind pairwise judge: both orders of every pair are sent, calibration items that are the same text twice
must come back as ties, and only the pairs on which both orders agree are counted. No judge verdict is committed yet, so none is reported.

Table~\ref{tab:probes} compares today's request (condition A, no context) with a request that adds the user memory (condition M) on a 4-billion-parameter model that runs on the
machine.\footnote{Both runs were made on an Apple M1 Pro with 32\,GB under macOS 27.0 with Ollama 0.34.2, as the manifest of each run records; the model has 3.9 billion parameters, quantised to Q4\_K\_M.} A second run of A gives the same 80 of 312. The memory moves the probes that only the memory can decide from $15.1\%$ to $55.6\%$ (a sign test on the
probes where the two runs disagree: 5 for A, 56 for M, $p<0.0001$) and leaves the controls and the probes decided by earlier lines unchanged
($p=0.81$). That is the shape the bench is built to show; it is not a claim about the in-place method, which is
indifferent to the engine, and the small model is a floor, not a recommendation.

\begin{table}[t]
\centering
\caption{Probes answered correctly, \texttt{gemma3:4b} on the machine, 312 probes. A: today's request. M: with the user memory.
Source: \texttt{runs/A-gemma3-4b} and \texttt{runs/M-gemma3-4b} of the translation bench.}\label{tab:probes}
\small
\begin{tabular}{@{}lrrr@{}}
\toprule
Probes decided by & n & A & M\\
\midrule
the memory only & 126 & 19 (15.1\%) & 70 (55.6\%)\\
earlier lines only & 133 & 26 (19.5\%) & 24 (18.0\%)\\
nothing (controls) & 53 & 35 (66.0\%) & 37 (69.8\%)\\
all & 312 & 80 (25.6\%) & 131 (42.0\%)\\
\bottomrule
\end{tabular}
\end{table}

\subsection{Engine speed}

The engine bench times the translation engines on four language pairs with 12 lines each, in batches of three. On the machine, Apple's low-latency
engine answers a batch of three lines in $0.15$\,s (warm median), and the cloud engines in $0.9$--$1.4$\,s, mostly network. It scores about
1.2 points below the cloud engines on a five-point scale of two model judges; differences under about $0.3$ are noise. That is the reason the scroll path of Section~\ref{sec:scroll}
does not wait for any engine: the cheapest engine is already slower than a frame.
